\label{sec:results}

The internship delivered a working, end-to-end ROS~2 pipeline that reconstructs a georeferenced 3D point cloud of the sea surface from a calibrated stereo pair and an independent navigation solution, and that extracts sea-state parameters from two independent sensing modalities. Against the objectives set out in Section~\ref{sec:objectives}: a working pipeline was built and demonstrated on real recorded data; six interchangeable disparity-estimation backends were implemented and systematically compared specifically to identify one that copes with the sea surface's low-texture, specular imaging conditions, which is the project's main technical contribution; and two independently derived, physically plausible wave estimates were obtained and cross-validated against each other without external ground truth, resolving along the way a series of unit, convention and noise-floor bugs that would otherwise have produced physically impossible results.

The pose/navigation estimate was likewise validated against an independent reference (the navigation filter's own modelled water level), with the resulting altitude residual converging to a small, stable value.

Relative to the original proposal, the modelling part of the work plan (state-of-the-art review, sensor characterization, estimation-algorithm development, validation) was completed and is the substance of this report; the wave-occurrence forecasting objective was descoped due to time-constraints in favour of consolidating reliable estimation first, and is carried forward as future work (Section~\ref{sec:future}).
